\documentclass[11pt,a4paper,twoside]{tesis}
% SI NO PENSAS IMPRIMIRLO EN FORMATO LIBRO PODES USAR
%\documentclass[11pt,a4paper]{tesis}

\usepackage{graphicx}
\usepackage[utf8]{inputenc}
\usepackage[spanish]{babel}
\usepackage[left=3cm,right=3cm,bottom=3.5cm,top=3.5cm]{geometry}

\begin{document}

%%%% CARATULA
% Comentar y descomentar según corresponda
\def\titulo{Licenciado }

\def\autor{Juan Manuel Pérez}
\def\tituloTesis{Mimetización entre interlocutores}
\def\runtitulo{Medición de la mimetización entre interlocutores utilizando series de tiempo}
\def\runtitle{Measuring entrainment between speakers using time series}
\def\director{Agustín Gravano}
\def\codirector{Ramiro Gálvez}
\def\lugar{Buenos Aires, 2015}
\newcommand{\HRule}{\rule{\linewidth}{0.2mm}}
%
\thispagestyle{empty}

\begin{center}\leavevmode

\vspace{-2cm}

\begin{tabular}{l}
\includegraphics[width=2.6cm]{logofcen.pdf}
\end{tabular}


{\large \sc Universidad de Buenos Aires

Facultad de Ciencias Exactas y Naturales

Departamento de Computaci\'on}

\vspace{6.0cm}

%\vspace{3.0cm}
%{
%\Large \color{red}
%\begin{tabular}{|p{2cm}cp{2cm}|}
%\hline
%& Pre-Final Version: \today &\\
%\hline
%\end{tabular}
%}
%\vspace{2.5cm}

{\huge\bf \tituloTesis}

\vspace{2cm}

{\large Tesis de \titulo en Ciencias de la Computaci\'on}

\vspace{2cm}

{\Large \autor}

\end{center}

\vfill

{\large

{Director: \director}

\vspace{.2cm}

{Codirector: \codirector}

\vspace{.2cm}

\lugar
}

% \newpage\thispagestyle{empty}


%%%% ABSTRACTS, AGRADECIMIENTOS Y DEDICATORIA
\frontmatter
\pagestyle{empty}
%\begin{center}
%\large \bf \runtitulo
%\end{center}
%\vspace{1cm}
\chapter*{\runtitulo}

\noindent El \emph{entrainment} (mimetización) es un fenómeno inconsciente que se manifiesta a través de la adaptación de posturas, forma de hablar, gestos faciales y otros comportamientos entre dos o más interactores. A su vez, la ocurrencia de esta mimetización está fuertemente emparentada con el sentimiento de empatía y compenetración entre los participantes.

En esta tesis, nos proponemos explorar una técnica algorítmica para detectar el entrainment entre variables prosódicas de dos personas. Esta técnica nos permitirá determinar si existe o no convergencia para ciertos parámetros, y ver como está esto correlacionado con variables sociales tales como la empatía, la compenetración con la tarea, y otras.

\bigskip

\noindent\textbf{Palabras claves:} Guerra, Rebelión, Wookie, Jedi, Fuerza, Imperio (no menos de 5).


\cleardoublepage
%\begin{center}
%\large \bf \runtitle
%\end{center}
%\vspace{1cm}
\chapter*{\runtitle}

\noindent In a galaxy far, far away, a psychopathic emperor and his most trusted servant -- a former Jedi Knight known as Darth Vader -- are ruling a universe with fear. They have built a horrifying weapon known as the Death Star, a giant battle station capable of annihilating a world in less than a second. When the Death Star's master plans are captured by the fledgling Rebel Alliance, Vader starts a pursuit of the ship carrying them. A young dissident Senator, Leia Organa, is aboard the ship \& puts the plans into a maintenance robot named R2-D2. Although she is captured, the Death Star plans cannot be found, as R2 \& his companion, a tall robot named C-3PO, have escaped to the desert world of Tatooine below. Through a series of mishaps, the robots end up in the hands of a farm boy named Luke Skywalker, who lives with his Uncle Owen \& Aunt Beru. Owen \& Beru are viciously murdered by the Empire's stormtroopers who are trying to recover the plans, and Luke \& the robots meet with former Jedi Knight Obi-Wan Kenobi to try to return the plans to Leia Organa's home, Alderaan. After contracting a pilot named Han Solo \& his Wookiee companion Chewbacca, they escape an Imperial blockade. But when they reach Alderaan's coordinates, they find it destroyed - by the Death Star. They soon find themselves caught in a tractor beam \& pulled into the Death Star. Although they rescue Leia Organa from the Death Star after a series of narrow escapes, Kenobi becomes one with the Force after being killed by his former pupil - Darth Vader. They reach the Alliance's base on Yavin's fourth moon, but the Imperials are in hot pursuit with the Death Star, and plan to annihilate the Rebel base. The Rebels must quickly find a way to eliminate the Death Star before it destroys them as it did Alderaan (aprox. 200 palabras).

\bigskip

\noindent\textbf{Keywords:} War, Rebellion, Wookie, Jedi, The Force, Empire (no menos de 5).

\cleardoublepage
\chapter*{Agradecimientos}

\noindent Lorem ipsum dolor sit amet, consectetur adipiscing elit. Fusce sapien ipsum, aliquet eget convallis at, adipiscing non odio. Donec porttitor tincidunt cursus. In tellus dui, varius sed scelerisque faucibus, sagittis non magna. Vestibulum ante ipsum primis in faucibus orci luctus et ultrices posuere cubilia Curae; Mauris et luctus justo. Class aptent taciti sociosqu ad litora torquent per conubia nostra, per inceptos himenaeos. Mauris sit amet purus massa, sed sodales justo. Mauris id mi sed orci porttitor dictum. Donec vitae mi non leo consectetur tempus vel et sapien. Curabitur enim quam, sollicitudin id iaculis id, congue euismod diam. Sed in eros nec urna lacinia porttitor ut vitae nulla. Ut mattis, erat et laoreet feugiat, lacus urna hendrerit nisi, at tincidunt dui justo at felis. Class aptent taciti sociosqu ad litora torquent per conubia nostra, per inceptos himenaeos. Ut iaculis euismod magna et consequat. Mauris eu augue in ipsum elementum dictum. Sed accumsan, velit vel vehicula dignissim, nibh tellus consequat metus, vel fringilla neque dolor in dolor. Aliquam ac justo ut lectus iaculis pharetra vitae sed turpis. Aliquam pulvinar lorem vel ipsum auctor et hendrerit nisl molestie. Donec id felis nec ante placerat vehicula. Sed lacus risus, aliquet vel facilisis eu, placerat vitae augue. % OPCIONAL: comentar si no se quiere

\cleardoublepage
\hfill \textit{A mis amigos.}

\hfill \textit{A mis compañeros de la Facultad, con los que tanto remamos.}

\hfill \textit{A mi familia, que me acompañaron.}  % OPCIONAL: comentar si no se quiere

\cleardoublepage
\tableofcontents

\mainmatter
\pagestyle{headings}

%%%% ACA VA EL CONTENIDO DE LA TESIS

\chapter{Introducción}

\section{Introducción}

Los sistemas de diálogo humano-computadora son cada vez más frecuentes, y sus aplicaciones comprenden una amplia gama de rubros: desde aplicaciones móviles, motores de búsqueda, juegos, o tecnologías de asistencia para ancianos y discapacitados.

Si bien es cierto que estos sistemas logran captar la dimensión lingüística de la comunicación humana, tienen un déficit importante a la hora de procesar y transmitir el aspecto superestructural de la comunicación, que radica en el intercambio de afecto, emociones, actitudes y otras intenciones de los participantes. La habilidad de los participantes de poder expresar, comprender, y reaccionar de acuerdo a estas señales sociales es necesaria para el entendimiento mutuo y una comunicación exitosa.

Un aspecto particular de la comunicación es el fenómeno de \emph{entrainment}(arrastre, mimetización, efecto camaleón), que comprende la adaptación inconsciente de las variables acústicas/prosódicas(a/p) (por ejemplo, el tono de la voz, la velocidad del habla, etc) de manera dinámica en el transcurso de una o varias interacciones. Este fenómeno ha sido introducido por \emph{Brennan et al}\cite{BRE1996} en 1996, y se ha observado que la convergencia de los participantes en estas variables ocurre en conjunto con una interacción más fluída y un mayor sentimiento de simpatía por sus interlocutores \cite{CHAR1999}.

Poder medir esta mimetización de los interlocutores no es una tarea fácil, sin embargo. En primer lugar, un diálogo no es una sucesión de turnos, sino que es una serie de tiempo dinámica, llena de interrupciones. Más aún, la mimetización no tiene un carácter instantáneo, sino que se sucede a lo largo de la interacción entre los participantes. Estos factores dificultan ostensiblemente poder modelar este fenómeno.

%%%% BIBLIOGRAFIA
\backmatter
\begin{thebibliography}{9}
\bibitem{BRE1996}
    Brennan:
    \emph{Lexical entrainment in spontaneous dialog},
    1996
\bibitem{CHAR1999}
    Chartrand, Bargh:
    \emph{The chameleon effect: The perception-behavior link and social interaction},
    1999
\bibitem{DEL2013}
    De Looze, Scherer, Vaughan, Campbell:
    \emph{Investigating automatic measurements of prosodic accommodation and its dynamics in social interaction},
    2013
\bibitem{KOU2008}
    Kousidis et al:
    \emph{Towards measuring continuous acoustic feature convergence in unconstrained spoken dialogues},
    2008
\bibitem{KOU2009}
    Kousidis, Dorran, McDonnell \& Coyle:
    \emph{Time Series Analysis of Acoustic Feature Convergence in Human Dialogues},
    2009
\bibitem{LEV2012}
    Levitan, Gravano, Willson, Benus, Hirschberg \& Nenkova,
    \emph{Acoustic-Prosodic entrainment and social behavior},
    2012
\bibitem{CHATFIELD}
    Chatfield C.,
    \emph{The analysis of time series: an introduction, Third Edition}
    1984
\end{thebibliography}
\end{document}
